% GENERATED FILE. Edit canonical lessons, then run npm run generate:books.
% canonical-source-hash: fnv1a64:0c834d2ed4e4b8de
% canonical-lessons: SA-C208-parivarah, SA-C208-bandhuh, SA-C208-sakha, SA-C208-satruh, SA-C208-prativesi

\chapter{Household, Relative, Companion, Enemy, Neighbour}
\label{ch:sa-household-relative-companion-enemy-neighbour}
\hlchaptermodality{208}

\begin{chapteropening}
\textbf{By the end of this chapter:} \emph{I can say a household, a relative, a companion, an enemy and a neighbour.}

Complete the last lesson of chapter 208: I can say a household, a relative, a companion, an enemy and a neighbour.
\end{chapteropening}

\section[\texorpdfstring{\sk{परिवारः} (parivāraḥ)}{परिवारः (parivāraḥ)}]{\sk{परिवारः} (parivāraḥ) — a household, a family}

\label{lesson:SA-C208-parivarah}



\begin{quote}
Before the new one: say the Sanskrit for advises, then the Sanskrit for honours.
\end{quote}

\subsection*{You'll want to know: \sk{परिवारः}}
\textbf{\sk{परिवारः}} — \emph{parivāraḥ} — \textquotedblleft{}a household, a family\textquotedblright{}.

\begin{grammarlens}[title={What you've built}]
One more word for this chapter.
\end{grammarlens}

\subsection*{Your turn}
\begin{itemize}
\raggedright
  \item \emph{Say it:} \emph{parivāraḥ}
  \item \emph{Say it:} \emph{parivāraḥ}, once more
  \item \emph{Say it:} say \emph{parivāraḥ}
  \item \emph{Recall:} say the Sanskrit for advises, then the Sanskrit for honours, then say \emph{parivāraḥ} again
\end{itemize}

\begin{tcolorbox}[breakable,colback=teal!4,colframe=teal!35!black,title={Before you move on}]
What does \sk{परिवारः} mean? (\textquotedblleft{}A household, a family\textquotedblright{}.) Say it once more.
\end{tcolorbox}
\section[\texorpdfstring{\sk{बन्धुः} (bandhuḥ)}{बन्धुः (bandhuḥ)}]{\sk{बन्धुः} (bandhuḥ) — a relative}

\label{lesson:SA-C208-bandhuh}



\begin{quote}
Before the new one: say the Sanskrit for honours, then the Sanskrit for a household.
\end{quote}

\subsection*{You'll want to know: \sk{बन्धुः}}
\textbf{\sk{बन्धुः}} — \emph{bandhuḥ} — \textquotedblleft{}a relative\textquotedblright{}.

\begin{grammarlens}[title={What you've built}]
One more word for this chapter.
\end{grammarlens}

\subsection*{Your turn}
\begin{itemize}
\raggedright
  \item \emph{Say it:} \emph{bandhuḥ}
  \item \emph{Say it:} \emph{bandhuḥ}, once more
  \item \emph{Say it:} say \emph{bandhuḥ}
  \item \emph{Recall:} say the Sanskrit for honours, then the Sanskrit for a household, then say \emph{bandhuḥ} again
\end{itemize}

\begin{tcolorbox}[breakable,colback=teal!4,colframe=teal!35!black,title={Before you move on}]
What does \sk{बन्धुः} mean? (\textquotedblleft{}A relative\textquotedblright{}.) Say it once more.
\end{tcolorbox}
\section[\texorpdfstring{\sk{सखा} (sakhā)}{सखा (sakhā)}]{\sk{सखा} (sakhā) — a companion, a (male) friend}

\label{lesson:SA-C208-sakha}



\begin{quote}
Before the new one: say the Sanskrit for a household, then the Sanskrit for a relative.
\end{quote}

\subsection*{You'll want to know: \sk{सखा}}
\textbf{\sk{सखा}} — \emph{sakhā} — \textquotedblleft{}a companion, a (male) friend\textquotedblright{}.

\begin{grammarlens}[title={What you've built}]
One more word for this chapter.
\end{grammarlens}

\subsection*{Your turn}
\begin{itemize}
\raggedright
  \item \emph{Say it:} \emph{sakhā}
  \item \emph{Say it:} \emph{sakhā}, once more
  \item \emph{Say it:} say \emph{sakhā}
  \item \emph{Recall:} say the Sanskrit for a household, then the Sanskrit for a relative, then say \emph{sakhā} again
\end{itemize}

\begin{tcolorbox}[breakable,colback=teal!4,colframe=teal!35!black,title={Before you move on}]
What does \sk{सखा} mean? (\textquotedblleft{}A companion, a (male) friend\textquotedblright{}.) Say it once more.
\end{tcolorbox}
\section[\texorpdfstring{\sk{शत्रुः} (śatruḥ)}{शत्रुः (śatruḥ)}]{\sk{शत्रुः} (śatruḥ) — an enemy}

\label{lesson:SA-C208-satruh}



\begin{quote}
Before the new one: say the Sanskrit for a relative, then the Sanskrit for a companion.
\end{quote}

\subsection*{You'll want to know: \sk{शत्रुः}}
\textbf{\sk{शत्रुः}} — \emph{śatruḥ} — \textquotedblleft{}an enemy\textquotedblright{}.

\begin{grammarlens}[title={What you've built}]
One more word for this chapter.
\end{grammarlens}

\subsection*{Your turn}
\begin{itemize}
\raggedright
  \item \emph{Say it:} \emph{śatruḥ}
  \item \emph{Say it:} \emph{śatruḥ}, once more
  \item \emph{Say it:} say \emph{śatruḥ}
  \item \emph{Recall:} say the Sanskrit for a relative, then the Sanskrit for a companion, then say \emph{śatruḥ} again
\end{itemize}

\begin{tcolorbox}[breakable,colback=teal!4,colframe=teal!35!black,title={Before you move on}]
What does \sk{शत्रुः} mean? (\textquotedblleft{}An enemy\textquotedblright{}.) Say it once more.
\end{tcolorbox}
\section[\texorpdfstring{\sk{प्रतिवेशी} (prativeśī)}{प्रतिवेशी (prativeśī)}]{\sk{प्रतिवेशी} (prativeśī) — a neighbour}

\label{lesson:SA-C208-prativesi}



\begin{quote}
Before the new one: say the Sanskrit for a companion, then the Sanskrit for an enemy.
\end{quote}

\subsection*{You'll want to know: \sk{प्रतिवेशी}}
\textbf{\sk{प्रतिवेशी}} — \emph{prativeśī} — \textquotedblleft{}a neighbour\textquotedblright{}.

\begin{grammarlens}[title={What you've built}]
One more word for this chapter.
\end{grammarlens}

\subsection*{Your turn}
\begin{itemize}
\raggedright
  \item \emph{Say it:} \emph{prativeśī}
  \item \emph{Say it:} \emph{prativeśī}, once more
  \item \emph{Say it:} say \emph{prativeśī}
  \item \emph{Recall:} say the Sanskrit for a companion, then the Sanskrit for an enemy, then say \emph{prativeśī} again
\end{itemize}

\begin{tcolorbox}[breakable,colback=teal!4,colframe=teal!35!black,title={Before you move on}]
What does \sk{प्रतिवेशी} mean? (\textquotedblleft{}A neighbour\textquotedblright{}.) Say it once more.
\end{tcolorbox}
